% Generated from validated Special post-draft synthesis. Do not hand-edit.
\section*{Monthly Signals}
\addcontentsline{toc}{section}{Monthly Signals}
\smallskip
\noindent \textbf{利用形態が広がった} 会話画面以外の利用形態が広がり、周辺の実装も重要になった。 \hfill{\footnotesize p.~\pageref{pkg:agent}}
\par
\smallskip
\noindent \textbf{モデル更新を系列で追う} 一般モデルと特化モデルの更新を別々の出来事として追う必要が強まった。 \hfill{\footnotesize p.~\pageref{pkg:model}}
\par
\smallskip
\noindent \textbf{配備の選択肢が広がった} 重み公開や提供形態の異なる選択肢が並立した。 \hfill{\footnotesize p.~\pageref{pkg:open}}
\par
\smallskip
\noindent \textbf{実行基盤も更新された} 推論実行基盤とライブラリも連続して更新された。 \hfill{\footnotesize p.~\pageref{pkg:runtime}}
\par
\smallskip
\noindent \textbf{評価と保護策も重視された} 機能拡大と並行して評価や保護策も重視された。 \hfill{\footnotesize p.~\pageref{pkg:control}}
\par
\medskip
\begin{claimboundary}[Retrospective scope]
本号は2025年後期を後日確認可能になった一次情報も用いて再構成するRetrospective Specialである。Coverage Periodと制作時点を同一視せず、vendor / project / author claimの境界は各記事で明示する。
\end{claimboundary}
\medskip
\tableofcontents
